\section{Related Work}
\label{sec:related}

Finite sampling can be addressed at three points in a quantum learning
pipeline. It can be addressed before execution, by deciding how to spend a
measurement budget. It can be addressed after averaging, by correcting the
statistics of expectation values. Or it can be addressed at the record
itself, by changing how repeated measurements become training data. Our
method lives at this third point, which is the least studied of the three.

The idea of averaging subsets of measurements is not new, and we do not claim
it. Cemin \emph{et al.} average subsets of measurements to enlarge a training
set for reduced quantum channel learning~\cite{cemin2024channels}, and it is
a classical result that training with noisy inputs behaves like an additional
penalty on large fitted weights~\cite{bishop1995training}. What has been
missing is the treatment of the aggregation scale as a tunable quantity. This
study evaluates the complete range from Raw to EV on one fixed QRC record,
calibrates the scale jointly with the readout regularization, gives an exact
account of the penalty this creates, and tests the result under three
allocations of a fixed execution budget on four matched forecasting tasks.
The contribution is the calibration and evaluation of the aggregation scale,
not the original introduction of subset averaging.

The other two points of the pipeline are well developed and complementary to
ours. Before execution, resource-aware optimization decides how many
measurements each part of a computation deserves~\cite{moussa2023resource}.
After averaging, Ahmed \emph{et al.} derive noise-aware training rules for
expectation values~\cite{ahmed2025optimal}, and Kreplin and Roth reduce
sampling noise through variance regularization~\cite{kreplin2024reduction}.
Fundamental limits of learning from finitely sampled physical systems have
also been mapped out~\cite{hu2023tackling}. None of these methods competes
with ours. Allocation ends where our method begins, and the averaged
quantities that noise-aware corrections operate on can themselves be produced
by grouped aggregation. Our method needs nothing beyond a single archived
measurement record.

Within QRC, a rich line of work studies how noise, dissipation, measurement,
finite Hilbert space, and input transformations shape temporal
features~\cite{kubota2022noise,palacios2024role,kubota2023temporal,martinezpena2023finite,mujal2023timeseries,suzuki2022natural,sannia2024dissipation,govia2022nonlinear,martinezpena2023capacity,kalfus2022hilbert,fry2023noiseinduced,govia2021singleoscillator}.
All of these studies analyze or modify the quantum representation. We do the
opposite. The representation is held completely fixed, and only the
organization of its repeated measurements changes. Because every reservoir in
those studies still produces a measurement record, adaptive shot aggregation
applies on top of any of them.
